%=====================================================================
\section{Final --- 21/07/2021}
%=====================================================================

% =====================================================================
\subsection{Ejercicio 1 --- Licitación transparente con AES-GCM (Tecnociudad)}
% =====================================================================

\begin{consigna}
El consejo deliberante del partido de Tecnociudad solicitó la construcción de un sistema
distribuido que permita realizar licitaciones de forma transparente y segura. Cuando el
partido necesita contratar cualquier tipo de servicios, genera un pliego con la descripción
de lo solicitado y espera que distintos contratistas presenten ofertas. Durante esta fase,
los costos son secretos para evitar que un oferente tenga preferencia. Una vez vencido el
plazo, se publica la propuesta de cada uno, junto con el precio, y se selecciona el que
cumple con las condiciones con el menor costo.

La implementación del sistema asegura que estas propiedades ocurran utilizando criptografía.
La descripción a alto nivel es la siguiente:
\begin{itemize}
  \item Se publica la licitación con un ID único asociado.
  \item Por cada licitación, se reciben archivos cifrados con AES-GCM de parte de los
  oferentes interesados. Los archivos incluyen un documento con el ID de la licitación, la
  descripción técnica, el CUIT de la empresa y el precio.
  \item Al terminar el periodo de la licitación los oferentes comparten la clave utilizada,
  y con ello se revisan todas las ofertas, y se adjudica la ganadora.
  \item Si un oferente no provee la clave, queda automáticamente descalificado.
  \item Para transparencia, se publican todas las ofertas recibidas ya descifradas y la
  adjudicada en un portal público que puede ser consultado libremente.
\end{itemize}

\begin{enumerate}
  \item[a)] Explicar por qué ningún oferente puede alterar su propuesta sin ser detectado.
  (1 pto.)
  \item[b)] Qué medidas adicionales serían necesarias para evitar que algún administrador
  del sistema ignore a un proveedor diciendo que nunca entregó la clave. (1 pto.)
  \item[c)] Considerar el caso de un oferente que envía un archivo a nombre de un competidor,
  con un precio irrisoriamente bajo que lo perjudicaría. \textquestiondown{}Cómo evitaría esto usando funciones
  criptográficas? (1 pto.)
\end{enumerate}
\end{consigna}

\paragraph{Temas.}
Primera Parte, Clase 3 --- MACs y Cifrado Autenticado (\S Modos AEAD en la práctica: GCM;
\S Definición de cifrado autenticado; \S No-repudio: por qué los MAC no lo dan). Primera
Parte, Clase 4 --- Asimétrico y Firma Digital (\S Diferencia entre cifrar y firmar;
\S RSA-Signature / firma del hash; \S PKI y certificados digitales, no-repudio).

\paragraph{Qué necesitás saber.}
\begin{itemize}
  \item \textbf{AES-GCM} es un modo AEAD (\emph{cifrado autenticado}): además de
  confidencialidad (modo Counter), produce una etiqueta de autenticación sobre el
  \textbf{texto cifrado}. Al descifrar, primero se verifica el tag; si el cifrado fue
  modificado, la verificación \textbf{falla} y el descifrado se rechaza (no se obtiene
  ``basura'' silenciosa como con cifrado no autenticado).
  \item \textbf{Encrypt-then-MAC / AEAD} da \textbf{confidencialidad + integridad +
  autenticación}, pero \textbf{con clave simétrica compartida NO da no-repudio}: quien
  verifica el tag/GHASH conoce la misma clave que quien lo generó, así que un tercero
  no puede distinguir quién produjo el mensaje (ambas partes podrían haberlo generado).
  \item Para \textbf{no-repudio} hace falta criptografía \textbf{asimétrica}: \textbf{firmar}
  (con la clave \textbf{privada} del firmante) es lo que da no-repudio, porque solo el dueño
  de la privada pudo producir esa firma, y cualquiera la verifica con la pública sin poder
  falsificarla.
\end{itemize}

\paragraph{Notas y conexiones.}
El escenario combina dos capas que el material distingue con cuidado: (i) AES-GCM protege
la \emph{confidencialidad e integridad del archivo cifrado en tránsito/reposo} durante la
licitación (nadie puede leer los precios antes de tiempo ni modificar un archivo ya
enviado sin que se note), pero (ii) por ser \textbf{simétrico y con clave compartida}, no
resuelve el problema de que un administrador (que también conoce/gestiona las claves)
niegue haber recibido algo, ni impide que alguien se haga pasar por otro oferente. Ambos
puntos requieren la capa de \textbf{firma digital / PKI} de la Clase 4, que sí da
autenticación de origen y no-repudio.

\paragraph{Resolución.}

\textbf{a)} Ningún oferente puede alterar su propuesta sin ser detectado porque el archivo
se envía cifrado con \textbf{AES-GCM}, que es un esquema de \textbf{cifrado autenticado}
(AEAD): además de cifrar el contenido (confidencialidad), GCM calcula una \textbf{etiqueta
de autenticación} sobre el texto cifrado (usando GHASH). Al momento de descifrar (una vez
compartida la clave), se \textbf{verifica primero el tag}; si el archivo cifrado fue
modificado en cualquier bit respecto del que se envió originalmente, la verificación del
tag \textbf{falla} y el descifrado se rechaza explícitamente (a diferencia de un cifrado
sin autenticar, donde la manipulación pasaría desapercibida y solo produciría datos
``basura'' al descifrar). Por lo tanto, cualquier intento de alterar el precio o la
descripción técnica \emph{después} de enviado el archivo es detectable al momento de
abrir las ofertas.

\textbf{b)} El esquema descripto usa una \textbf{clave simétrica compartida} (el oferente
``comparte la clave utilizada''); como se remarca en la Clase 3, un esquema con clave
compartida (MAC/AEAD simétrico) \textbf{no da no-repudio}: quien controla el sistema
(el administrador) también podría --- al menos en principio --- estar en condiciones de
negar la recepción de la clave sin que el oferente pueda probar lo contrario, ya que no hay
una firma verificable por un tercero independiente. La medida adicional es introducir
\textbf{firma digital (asimétrica)}: que el oferente, al enviar la clave de descifrado,
la \textbf{firme con su clave privada}, y que el administrador, al recibirla, devuelva un
\textbf{acuse de recibo firmado} con su propia clave privada. De esta forma:
\begin{itemize}
  \item La firma del oferente sobre el envío de la clave prueba, ante cualquier tercero
  (verificando con la clave pública del oferente), que efectivamente la envió.
  \item El acuse de recibo firmado por el administrador (con su clave privada, verificable
  con su pública conocida vía PKI/certificado) es una prueba de que el sistema la recibió;
  como solo el administrador pudo generar esa firma, \textbf{no puede negarlo después}
  (no-repudio).
\end{itemize}
Esto traslada el problema de ``palabra contra palabra'' a un esquema con \textbf{propiedades
de no-repudio}, que el cifrado simétrico por sí solo no ofrece.

\textbf{c)} El problema es que el sistema, tal como está descripto, no verifica que quien
envía un archivo ``a nombre de'' un CUIT sea realmente el dueño de esa identidad: cualquiera
podría escribir el CUIT de un competidor dentro del documento cifrado. Esto se evita
usando \textbf{firma digital}: cada oferente debe \textbf{firmar su propuesta} (el
documento completo: ID de licitación, descripción técnica, CUIT y precio) con su
\textbf{clave privada} antes de cifrarlo y enviarlo, y esa clave privada debe estar
asociada a su identidad mediante un \textbf{certificado digital} emitido por una
\textbf{CA} (PKI), tal como en el ejemplo de firma de binarios de la Clase 4. Al
adjudicar, el administrador (o cualquiera con la clave pública certificada del oferente)
\textbf{verifica la firma}: si el documento dice ser del competidor pero está firmado con
la clave privada del atacante, la verificación con la clave pública \emph{del competidor}
\textbf{falla}, y la propuesta fraudulenta se descarta. En síntesis: \textbf{AES-GCM} da
confidencialidad e integridad del archivo, pero la \textbf{autenticidad del emisor} (que
el CUIT declarado sea quien realmente presentó la oferta) requiere \textbf{firma digital
con PKI}, que es justamente lo que impide la suplantación de identidad descripta.

% =====================================================================
\subsection{Ejercicio 2 --- Elección de supervisor mediante hash y XOR}
% =====================================================================

\begin{consigna}
El siguiente protocolo fue propuesto para que $N$ nodos de un sistema puedan elegir un
supervisor, de tal manera que ningún nodo pueda ganar o perder sistemáticamente
manipulando el intercambio:
\begin{enumerate}
  \item Cada nodo genera un valor $r_i$ de 64 bits y envía al resto $H(r_i)$.
  \item Luego de recibir todos los valores, cada nodo publica $r_i$.
  \item Cada nodo calcula $r_x = r_1 \oplus r_2 \oplus \dots \oplus r_n$.
\end{enumerate}
El nodo a menor distancia de $r_x$ es el nuevo supervisor, utilizando la distancia de
Hamming como métrica (se comparan bit a bit los dígitos binarios de los dos valores, y
cada discrepancia suma 1 de distancia --- por ejemplo 0010 y 1110 están a distancia 2).
En caso de empate, se reinicia.

\begin{enumerate}
  \item[a)] Explicar por qué el paso 1 es necesario y no se puede comenzar el intercambio
  en el paso 2 directamente. (1 pto.)
  \item[b)] Considerar el caso de solo tres nodos. \textquestiondown{}Qué pasaría si el nodo 2, por ejemplo,
  decide abusar el sistema, y repite en el paso 1 el valor $H(r_1)$ como si fuese el suyo?
  Proponer una modificación simple que evite el problema. (1,5 ptos.)
  \item[c)] Proponer una mejora al algoritmo para que no sea necesario repetir todo el
  intercambio en caso de empate y siempre se llegue a un resultado que no pueda ser
  manipulado. (0.5 ptos.)
\end{enumerate}
\end{consigna}

\paragraph{Temas.}
Primera Parte, Clase 3 --- MACs y Cifrado Autenticado (\S Funciones de hash: resistencia a
preimagen, segunda preimagen y colisiones; uso como \emph{compromiso}/commitment previo a
revelar un valor). \textbf{Nota:} el protocolo en sí (elección de supervisor con XOR y
distancia de Hamming) no está descripto literalmente en el material como esquema con
nombre propio; la resolución se apoya en las \textbf{propiedades de las funciones de hash}
que sí están desarrolladas en la Clase 3, aplicadas como un esquema de compromiso
(committing a $r_i$ vía $H(r_i)$ antes de revelarlo).

\paragraph{Qué necesitás saber.}
\begin{itemize}
  \item \textbf{Resistencia a preimagen:} dado un hash $y$, es computacionalmente inviable
  hallar un mensaje $x$ tal que $H(x)=y$. Este es el uso relevante acá: publicar $H(r_i)$
  primero \emph{compromete} al nodo con un valor sin revelarlo, porque nadie puede, a partir
  del hash, deducir $r_i$ para ajustar su propio valor en función de los demás.
  \item \textbf{Resistencia a segunda preimagen:} dado un mensaje $x$, es inviable hallar
  $x'\neq x$ con el mismo hash. Esto es lo que impide, una vez publicado $H(r_i)$, cambiar
  el $r_i$ real por otro $r_i'$ que también hashee a lo mismo en el paso 2.
  \item El material señala explícitamente el uso de compromisos criptográficos (``probar que
  se tenía algo antes de revelarlo''; ``juegos de cartas con compromiso'') como aplicación de
  la resistencia a (segunda) preimagen.
\end{itemize}

\paragraph{Notas y conexiones.}
El paso 1 (publicar $H(r_i)$ antes de $r_i$) es exactamente el patrón de
\textbf{compromiso-y-revelación} (commit-then-reveal) que la Clase 3 asocia a la resistencia
a preimagen/segunda preimagen: sirve para que ningún participante pueda elegir su valor
\emph{después} de ver los de los demás (lo cual rompería la imparcialidad de la elección).
Es el mismo principio que sustenta, por ejemplo, no poder generar de antemano el $r_i$ en
función del resultado final.

\paragraph{Resolución.}

\textbf{a)} El paso 1 (publicar $H(r_i)$ antes de $r_i$) es necesario porque, si en cambio
se comenzara directamente publicando los $r_i$ en claro (paso 2), un nodo malicioso podría
esperar a ver los valores $r_j$ ($j\neq i$) que ya publicaron todos los demás y \textbf{recién
entonces} elegir su propio $r_i$ de forma tal que el XOR resultante $r_x$ quede
convenientemente cerca (en distancia de Hamming) de su propio valor, garantizándose ser
supervisor. Publicar primero $H(r_i)$ actúa como un \textbf{compromiso}: por
\textbf{resistencia a preimagen}, ningún otro nodo puede, a partir de $H(r_i)$, obtener
$r_i$ antes de que se revele en el paso 2, de modo que nadie puede ajustar su propio valor
en función de los $r_i$ de los demás. Así se garantiza que todos los $r_i$ quedaron
``congelados'' antes de que ninguno conociera el de los otros, haciendo el intercambio
imparcial.

\textbf{b)} Con tres nodos, si el nodo 2 en el paso 1 publica $H(z)=H(r_1)$ (repite el
compromiso del nodo 1) en lugar de comprometerse con su propio valor, en el paso 2 puede
publicar $z=r_1$ como si fuera su $r_2$ (es un valor válido porque su hash coincide con lo
publicado). El resultado sería $r_x = r_1\oplus r_1\oplus r_3 = r_3$ (el XOR de $r_1$ consigo
mismo se cancela), y el nodo 2 conocería de antemano que $r_x=r_3$: puede entonces mirar la
distancia de Hamming entre $r_1,r_2^{\text{real}}$ y $r_3$ contra $r_x=r_3$ y decidir
manipular el resultado a su favor (por ejemplo, el propio nodo 3 quedaría siempre exactamente
a distancia 0 de $r_x$, y el nodo 2 puede jugar con esa certeza, o directamente evitar que el
nodo 1 o 3 ganen si le conviniera). En cualquier caso, se rompe la propiedad buscada de que
``ningún nodo pueda ganar o perder sistemáticamente manipulando el intercambio''.

\textbf{Modificación simple:} exigir que \textbf{todos los compromisos $H(r_i)$ publicados en
el paso 1 sean distintos entre sí} ($H(r_i)\neq H(r_j)$ para todo $i\neq j$); si dos nodos
publican el mismo hash, se detecta la colisión de compromisos y se reinicia el proceso (o se
descalifica al nodo repetido). Como por \textbf{resistencia a segunda preimagen} es inviable
para el nodo 2 encontrar un $r_2'\neq r_1$ propio con $H(r_2')=H(r_1)$, la única forma de
publicar un hash igual al de otro nodo es copiándolo literalmente, lo cual esta verificación
detecta y bloquea antes de llegar al paso 2.

\textbf{c)} Para no tener que reiniciar todo el intercambio (generar nuevos $r_i$ y repetir
los tres pasos) en caso de empate en la distancia de Hamming, se puede definir un
\textbf{criterio de desempate determinístico y no manipulable} sobre la información que ya
está publicada, sin generar nueva aleatoriedad: por ejemplo, entre los nodos empatados en
distancia mínima a $r_x$, elegir como supervisor al que tenga \textbf{el valor $r_i$
numéricamente mayor (o menor)} interpretado como entero, o el que primero coincide bit a bit
con $r_x$ leyendo de izquierda a derecha. Como los $r_i$ ya fueron comprometidos con hash y
revelados (y no pueden cambiarse sin romper la resistencia a segunda preimagen del
compromiso), este criterio de desempate no puede ser manipulado por ningún nodo a esta
altura del protocolo (todos los valores ya están fijados e inmutables), y siempre produce un
único ganador sin necesidad de reiniciar el intercambio completo.

% =====================================================================
\subsection{Ejercicio 3 --- Pentesting con hipótesis de falla (Agrolin)}
% =====================================================================

\begin{consigna}
Agrolin S.A.\ se dedica a la construcción de soluciones tecnológicas para el campo.

Recientemente se publicó un informe que menciona múltiples vulnerabilidades en uno de sus
productos y provocó muchas pérdidas de ventas a punto de concretarse. Para corregir la
situación, la empresa decidió contratar un servicio de \emph{pen testing}, realizar una
evaluación más exhaustiva y corregir el producto para empezar a recuperar el daño a su
imagen.

El producto en cuestión permite hacer un seguimiento del rendimiento de las cosechas a lo
largo del tiempo y determinar las rotaciones óptimas para mejorar la productividad. El
sistema está compuesto por un frontend web que es un \emph{single page web app}, donde el
usuario puede registrarse, autenticarse, dar de alta lotes marcado sus límites en un mapa,
actualizarlos o eliminarlos. El sistema muestra según datos históricos de todos los clientes
de la zona, el rinde esperado para la superficie en cuestión, y sugiere el tipo de cultivo
con mayor rentabilidad en función de los valores de rinde y precios de mercados futuros
estimados, y permite que el usuario ingrese los datos del tipo de cultivo y luego de la
cosecha el rinde obtenido para comparar y poder mejorar en las próximas sugerencias.

La aplicación web se ejecuta en los navegadores de cada cliente, y hace llamadas a un
servidor backend a través de una API REST.

Siguiendo la metodología de hipótesis de falla, establecer 4 hipótesis de vulnerabilidades
que tengan alta probabilidad de existir en el escenario descripto. Para cada una, describir
únicamente la hipótesis y la prueba que realizaría para confirmarla o refutarla (1 pto.\ por
hipótesis).

\emph{Aclaraciones:} No se considerarán válidas hipótesis genéricas. Tampoco se considerarán
dos hipótesis que sean variantes del mismo tipo de falla (por ejemplo, el mismo problema en
dos pantallas distintas).
\end{consigna}

\paragraph{Temas.}
Segunda Parte, Clase 6 --- Pentesting (\S Metodología de Hipótesis de Falla: los pasos;
\S Paso 2 --- Planteo de hipótesis de falla). Segunda Parte, Clase 3 --- Principios de
Diseño y Vulnerabilidades (\S CWE / Top 25 MITRE: SQL Injection CWE-89; \S OWASP Top 10;
\S STRIDE: Denial of Service, Spoofing/Improper Authentication).

\paragraph{Qué necesitás saber.}
\begin{itemize}
  \item La \textbf{metodología de hipótesis de falla} exige, en el paso de
  \textbf{planteo de hipótesis}, asumir la existencia de una vulnerabilidad concreta
  \emph{antes} de probarla; luego se define una \textbf{prueba} específica para
  confirmarla o refutarla (no alcanza con nombrar la categoría genérica: hay que decir
  \emph{qué} se sospecha y \emph{cómo} se lo comprueba en este sistema puntual).
  \item El material lista, dentro del \textbf{Top 25 CWE de MITRE}, entre las debilidades
  más comunes: \textbf{CWE-89 SQL Injection} (falta de control en el ingreso de datos que
  permite ejecutar comandos SQL: los datos del usuario se mezclan sin sanitizar con el
  comando y la base los interpreta como código; defensa: \emph{prepared statements}/
  parametrizar), \textbf{CWE-79 Cross-site Scripting}, y (Clase 3, taxonomía STRIDE)
  \textbf{Denial of Service} (agotar los recursos necesarios para dar el servicio,
  afectando disponibilidad) como categoría de amenaza reconocida explícitamente como
  \emph{ataque de seguridad} aunque no comprometa confidencialidad/integridad.
  \item \textbf{OWASP} publica el Top 10 de vulnerabilidades web más comunes (frameworks de
  referencia junto con MITRE), relevante para un frontend \emph{single page web app} +
  API REST como el descripto (autenticación, autorización, control de acceso a recursos por
  ID).
\end{itemize}

\paragraph{Notas y conexiones.}
El enunciado da todas las pistas típicas de una aplicación web con API REST y autenticación
de usuarios: (1) el sistema tiene \textbf{autenticación} (registrarse/autenticarse), (2)
tiene \textbf{ingreso de datos de usuario} que probablemente lleguen a una base de datos
(lotes, límites en mapa, tipo de cultivo, rinde), (3) hace \textbf{consultas agregadas
costosas} sobre datos históricos de todos los clientes de la zona, y (4) maneja
\textbf{recursos por identificador} (lotes) que pertenecen a distintos clientes. Cada una
de estas pistas mapea a una categoría de vulnerabilidad del Top 25/CWE o de STRIDE vista en
la Clase 3, y la consigna pide aplicar sobre ellas el paso de \textbf{hipótesis + prueba}
de la Clase 6 (no basta con nombrar la categoría).

\paragraph{Resolución.}

\textbf{Hipótesis 1 --- Autenticación deficiente (Improper Authentication).}
\emph{Hipótesis:} el sistema de autenticación del frontend/API es débil o directamente
inexistente en algunos endpoints, permitiendo iniciar sesión con credenciales triviales o
acceder a funciones que deberían requerir sesión válida sin ella.
\emph{Prueba:} intentar registrarse con credenciales simples (usuario/contraseña
predecibles) y verificar si el sistema las acepta sin exigir complejidad; adicionalmente,
intentar invocar directamente los endpoints de la API REST (alta/consulta de lotes) sin
enviar token/cookie de sesión, para verificar si el backend responde igual que a un usuario
autenticado.

\textbf{Hipótesis 2 --- SQL Injection (CWE-89) en el ingreso de datos.}
\emph{Hipótesis:} los campos donde el usuario ingresa datos (tipo de cultivo, límites del
lote, rinde obtenido) no sanitizan/parametrizan la entrada antes de construir la consulta
a la base de datos, permitiendo inyectar comandos SQL.
\emph{Prueba:} en un campo de texto (por ejemplo, nombre del cultivo), ingresar una cadena
con sintaxis SQL (p.\ ej.\ un valor que cierre la comilla e intente insertar/alterar una
tabla) y verificar, sin necesidad de dañar la base, si el sistema devuelve un error de
sintaxis SQL expuesto o un comportamiento anómalo que indique que el dato fue interpretado
como comando en lugar de como valor.

\textbf{Hipótesis 3 --- Denial of Service por consultas agregadas costosas.}
\emph{Hipótesis:} la funcionalidad que muestra, según datos históricos de \textbf{todos los
clientes de la zona}, el rinde esperado y sugiere el cultivo más rentable en función de
precios de mercados futuros, requiere procesar/parsear múltiples consultas pesadas por cada
solicitud; enviar varias de estas solicitudes en paralelo podría agotar los recursos del
servidor y denegar el servicio a otros usuarios.
\emph{Prueba (acotada, sin ejecutar el ataque real):} verificar únicamente si es posible
disparar 2 de estas consultas costosas en paralelo desde una misma sesión/IP sin que el
sistema las limite (rate limiting) o las rechace; no se realiza el ataque a mayor escala
para no afectar el servicio, y se reporta la situación si se confirma la ausencia de
límites.

\textbf{Hipótesis 4 --- Autorización deficiente / control de acceso a recursos de otro
cliente (Improper Authorization / Broken Access Control).}
\emph{Hipótesis:} una vez autenticado, el sistema no verifica correctamente que el usuario
solo pueda acceder a sus propios lotes, permitiendo acceder a lotes (y sus límites en el
mapa, rindes, etc.) de otros clientes modificando un identificador en la URL o en la
llamada a la API REST.
\emph{Prueba:} tras crear una cuenta y autenticarse, identificar la URL/endpoint que expone
un lote propio con su identificador (p.\ ej.\ \texttt{.../cliente/12/campo/1}) y probar
acceder al mismo recurso cambiando el identificador de cliente por uno ajeno (p.\ ej.\
\texttt{.../cliente/1/campo/1}); si el sistema devuelve los datos sin validar que el lote
pertenece al usuario autenticado, se confirma la hipótesis. No se modifican datos ajenos
para no ser intrusivo; solo se confirma el acceso indebido.

\textbf{Nota sobre alcance:} las cuatro hipótesis corresponden a categorías distintas de
falla (autenticación, inyección de datos, disponibilidad y autorización), tal como exige la
aclaración de la consigna, y todas están sustentadas en la taxonomía CWE/STRIDE/OWASP
desarrollada en la Clase 3 y en el paso de ``planteo de hipótesis'' de la metodología de
pentesting de la Clase 6.
